\documentclass[runningheads]{llncs}

\usepackage[review,year=2026,ID=*****]{eccv}

\usepackage{eccvabbrv}
\usepackage{graphicx}
\usepackage{booktabs}
\usepackage{amsmath,amssymb}
\usepackage[accsupp]{axessibility}
\usepackage[pagebackref,breaklinks,colorlinks,citecolor=eccvblue]{hyperref}

\newcommand{\modelname}{\textsc{FiHaRD}}

\title{Supplementary Material for \modelname}
\titlerunning{Supplementary Material for \modelname}

\author{Anonymous ECCV/ACVR Submission}
\authorrunning{Anonymous ECCV/ACVR Submission}
\institute{Paper ID *****}

\begin{document}

\maketitle

\section{Denoiser Architecture}
\label{sec:supp_denoiser}

This section documents the denoiser used in the diffusion refinement stage of \modelname{}. The network takes the noised residual state in spatio-spectral coordinates together with conditioning derived from the coarse prediction, and estimates the corresponding $v$-target. Its design separates spatial interaction across joints from temporal evolution across the forecast horizon, while preserving explicit joint identity throughout the network.

\subsection{Input Tokenization}

Let
\[
z_t \in \mathbb{R}^{B \times T_{\mathrm{out}} \times F},
\qquad
F = 3N,
\]
denote the noised residual state at diffusion timestep $t$. We reshape it into joint-coordinate triplets,
\[
\tilde{z}_t \in \mathbb{R}^{B \times T_{\mathrm{out}} \times N \times 3},
\]
and project each triplet independently into a latent space of dimension $d$:
\[
X^{(0)} = W_{\mathrm{in}} \tilde{z}_t + P_{\mathrm{time}} + P_{\mathrm{joint}},
\]
where $W_{\mathrm{in}} \in \mathbb{R}^{3 \times d}$ is a learned input projection, $P_{\mathrm{time}} \in \mathbb{R}^{1 \times T_{\mathrm{out}} \times 1 \times d}$ is a learned embedding over the prediction horizon, and $P_{\mathrm{joint}} \in \mathbb{R}^{1 \times 1 \times N \times d}$ is a learned joint embedding. The resulting tensor
\[
X^{(0)} \in \mathbb{R}^{B \times T_{\mathrm{out}} \times N \times d}
\]
contains one latent token per joint and per future timestep.

\subsection{Diffusion Timestep and Coarse-Future Conditioning}

The diffusion timestep is encoded by a sinusoidal embedding followed by a two-layer MLP:
\[
e_t = \mathrm{MLP}(\mathrm{SinEmbed}(t)),
\qquad
e_t \in \mathbb{R}^{B \times d}.
\]
This embedding modulates every denoiser block through feature-wise affine conditioning. For a hidden tensor
\[
X \in \mathbb{R}^{B \times T_{\mathrm{out}} \times N \times d},
\]
the modulation is
\[
\mathrm{FiLM}(X,e_t)
=
X \odot \bigl(1+\Gamma(e_t)\bigr) + \beta(e_t),
\]
with learned parameters $\Gamma(e_t), \beta(e_t) \in \mathbb{R}^{B \times 1 \times 1 \times d}$ broadcast over time and joints.

The coarse forecast is projected into a conditioning memory
\[
M \in \mathbb{R}^{B \times T_{\mathrm{out}} \times N \times d},
\]
which is used by cross-attention. Because this memory preserves both temporal index and joint identity, refinement remains anchored to the coarse future instead of replacing it with an unconstrained latent sample.

\subsection{Factorized Denoiser Blocks}

Each denoiser block applies three attention stages in sequence:
\[
\begin{aligned}
Y_{\mathrm{sp}}
&=
\mathrm{SpAttn}\!\left(
\mathrm{FiLM}_{\mathrm{sp}}(X^{(\ell-1)},e_t)
\right),\\
Y_{\mathrm{tmp}}
&=
\mathrm{TempAttn}\!\left(
\mathrm{FiLM}_{\mathrm{tmp}}(Y_{\mathrm{sp}},e_t)
\right),\\
X^{(\ell)}
&=
\mathrm{CrossAttn}\!\left(
\mathrm{FiLM}_{\mathrm{cr}}(Y_{\mathrm{tmp}},e_t), M
\right).
\end{aligned}
\]
Different FiLM projections are used before the spatial, temporal, and cross-attention stages. The block is therefore factorized rather than using full joint-time attention: spatial self-attention mixes joints at fixed future timesteps, temporal self-attention mixes future timesteps for each joint independently, and cross-attention injects coarse-future memory while preserving joint identity.

\subsection{Attention Pattern}

\paragraph{Spatial self-attention.}
For a fixed future timestep $\tau$, let
\[
X_\tau \in \mathbb{R}^{N \times d}
\]
denote the matrix of joint tokens. Spatial attention is applied across the joint dimension using standard multi-head scaled dot-product attention. This stage captures instantaneous inter-joint dependencies without mixing temporal information.

\paragraph{Temporal self-attention.}
For a fixed joint $j$, let
\[
X_j \in \mathbb{R}^{T_{\mathrm{out}} \times d}
\]
denote its latent future trajectory. Temporal attention is then applied across the prediction horizon. This stage models how each articulator evolves over time while keeping joint identities explicit.

\paragraph{Temporal cross-attention.}
After the two self-attention stages, conditioning is injected through temporal cross-attention. For each joint $j$, the query sequence attends only to the memory sequence of the same joint:
\[
Q_{j}^{\mathrm{cr}} = \tilde{X}_j W_Q^{\mathrm{cr}},
\qquad
K_{j}^{\mathrm{cr}} = M_j W_K^{\mathrm{cr}},
\qquad
V_{j}^{\mathrm{cr}} = M_j W_V^{\mathrm{cr}}.
\]
This preserves anatomical correspondence by refining each predicted joint trajectory from the temporal memory of the matching coarse joint trajectory, rather than mixing conditioning information across different articulators.

\subsection{Residual Structure and Output Projection}

In implementation, each attention stage is wrapped in a Transformer-style residual block with normalization and a position-wise MLP:
\[
X \leftarrow X + \mathrm{Attn}\bigl(\mathrm{Norm}(\mathrm{FiLM}(X,e_t))\bigr),
\]
followed by
\[
X \leftarrow X + \mathrm{MLP}\bigl(\mathrm{Norm}(X)\bigr).
\]
This pattern is repeated for the spatial, temporal, and cross-attention sublayers.

After the final block, the hidden tensor
\[
X^{(L)} \in \mathbb{R}^{B \times T_{\mathrm{out}} \times N \times d}
\]
is projected back to joint coordinates:
\[
\hat{v} = W_{\mathrm{out}} X^{(L)},
\qquad
W_{\mathrm{out}} \in \mathbb{R}^{d \times 3}.
\]
Flattening the joint and coordinate dimensions yields the final prediction
\[
\hat{v} \in \mathbb{R}^{B \times T_{\mathrm{out}} \times F},
\]
which is trained against the $v$-target defined by the spatio-spectral forward process. In the experiments reported in the main paper, the denoiser uses latent width $512$, depth $6$, and $4$ attention heads.

\section{Theoretical Properties of the Spatio-Spectral Forward Process}

\begin{proposition}[Validity of the spatio-spectral forward process]
\label{prop:ss_validity}
Let
\[
\Sigma_{\mathrm{space}} = U \Lambda U^\top,
\qquad
\Lambda = \operatorname{diag}(\lambda_1,\dots,\lambda_F),
\]
with $U$ orthonormal and $\lambda_m \geq 0$, and let $\nu_\ell \geq 0$ denote the frequency weights associated with the future DCT modes. Define the joint mode energies
\[
\kappa_{\ell,m} = \nu_\ell \lambda_m,
\]
their mean
\[
\bar{\kappa}
=
\frac{1}{T_{\mathrm{out}}F}\sum_{\ell,m}\kappa_{\ell,m},
\]
and the mode-dependent rates
\[
\rho_{\ell,m}
=
\left(\frac{\kappa_{\ell,m}}{\bar{\kappa}}\right)^p,
\qquad p>0.
\]
Given a base cumulative schedule $0<\bar{\alpha}_t\leq 1$, define
\[
\bar{\alpha}_{t,\ell,m} = \bar{\alpha}_t^{\,\rho_{\ell,m}},
\]
and assume that the cumulative retention factors satisfy $0<\bar{\alpha}_{t,\ell,m}\leq 1$ for every timestep $t$ and every spatio-spectral mode $(\ell,m)$. Then the forward process
\[
z_t
=
\sqrt{\bar{\alpha}_t}\odot z_0
+
\sqrt{1-\bar{\alpha}_t}\odot \epsilon,
\qquad
\epsilon \sim \mathcal{N}(0,I),
\]
with $z_0=r_0U$, is:
\begin{enumerate}
    \item admissible, since each mode has nonnegative variance $1-\bar{\alpha}_{t,\ell,m}$;
    \item linear-Gaussian in the residual coordinates, because it is diagonal Gaussian in $z$-space and mapped back by the orthonormal transform $r=zU^\top$;
    \item Markov, with one-step coefficients $\alpha_{t,\ell,m}=\bar{\alpha}_{t,\ell,m}/\bar{\alpha}_{t-1,\ell,m}$ defined mode by mode; and
    \item compatible with $v$-prediction, since together with $z_t$, the target
    \[
    v_t
    =
    \sqrt{\bar{\alpha}_t}\odot \epsilon
    -
    \sqrt{1-\bar{\alpha}_t}\odot z_0
    \]
    uniquely determines both $z_0$ and $\epsilon$ mode by mode.
\end{enumerate}
\end{proposition}

\section{Proof of Proposition~\ref{prop:ss_validity}}

We prove each statement in turn.

\paragraph{Admissibility.}
By assumption, every cumulative retention factor satisfies
\[
0<\bar{\alpha}_{t,\ell,m}\leq 1.
\]
Therefore each mode of the forward process has variance
\[
1-\bar{\alpha}_{t,\ell,m}\geq 0,
\]
so the noising rule
\[
z_{t,\ell,m}
=
\sqrt{\bar{\alpha}_{t,\ell,m}}\,z_{0,\ell,m}
+
\sqrt{1-\bar{\alpha}_{t,\ell,m}}\,\epsilon_{\ell,m}
\]
defines a valid Gaussian perturbation for every spatio-spectral mode.

\paragraph{Linear-Gaussian structure.}
In the spectral coordinates, the process is diagonal across modes and Gaussian because $\epsilon\sim\mathcal{N}(0,I)$ and the update is affine in $(z_0,\epsilon)$. The residual coordinates are obtained through the linear map
\[
r_t = z_t U^\top,
\]
with $U$ orthonormal. Linear images of Gaussian variables remain Gaussian, hence the induced process in residual coordinates is also linear-Gaussian.

\paragraph{Markov property.}
Let
\[
\alpha_{t,\ell,m}
=
\frac{\bar{\alpha}_{t,\ell,m}}{\bar{\alpha}_{t-1,\ell,m}}.
\]
Then the cumulative coefficients factorize exactly as in standard diffusion, mode by mode. Consequently, each coordinate admits a one-step transition of the form
\[
z_{t,\ell,m}
=
\sqrt{\alpha_{t,\ell,m}}\,z_{t-1,\ell,m}
+
\sqrt{1-\alpha_{t,\ell,m}}\,\eta_{t,\ell,m},
\qquad
\eta_{t,\ell,m}\sim\mathcal{N}(0,1),
\]
which depends on the past only through $z_{t-1,\ell,m}$. Since this holds independently for every mode, the full process is Markov.

\paragraph{Compatibility with $v$-prediction.}
For each mode, define
\[
A_{t,\ell,m}=\sqrt{\bar{\alpha}_{t,\ell,m}},
\qquad
B_{t,\ell,m}=\sqrt{1-\bar{\alpha}_{t,\ell,m}}.
\]
The forward process and the $v$-target are
\[
z_{t,\ell,m}=A_{t,\ell,m}z_{0,\ell,m}+B_{t,\ell,m}\epsilon_{\ell,m},
\]
\[
v_{t,\ell,m}=A_{t,\ell,m}\epsilon_{\ell,m}-B_{t,\ell,m}z_{0,\ell,m}.
\]
This is a $2\times2$ linear system in the unknowns $(z_{0,\ell,m},\epsilon_{\ell,m})$. Its coefficient matrix
\[
\begin{bmatrix}
A_{t,\ell,m} & B_{t,\ell,m}\\
-B_{t,\ell,m} & A_{t,\ell,m}
\end{bmatrix}
\]
has determinant
\[
A_{t,\ell,m}^2+B_{t,\ell,m}^2
=
\bar{\alpha}_{t,\ell,m} + (1-\bar{\alpha}_{t,\ell,m})
=
1.
\]
Hence it is invertible, so $(z_{0,\ell,m},\epsilon_{\ell,m})$ is uniquely determined by $(z_{t,\ell,m},v_{t,\ell,m})$ for every mode. Therefore the proposed forward process is compatible with standard $v$-prediction.

\end{document}
